\scanpage{69}
\uline{Proposition} Soit $D$ un diviseur satisfaisant a) et b), et soit $M$ la famille des \uline{diviseurs $>0$} \uline{numériquement} équivalents à $D$. Alors tout diviseur \unsure{numériqt|numériquement} équivalent à $0$ est lin. équivalent à un diviseur de la forme $D'-D$, avec $D'\in M$, \deleted{\unsure{donc}} \unsure{i.e|i.e.} $M$ est une classe mod $\underline{\mathrm{Pic}}^n_{X/k}$.

Le th. de \deleted{finitude} R.R. \inserted{(faible) sur $X$} montre que les polynômes de Hilbert des $X\in M$ sont les \unsure{ms\%}. Le th. de finitude de Chow (repris ds les schémas de Hilbert) montre donc que $M$ est quasi-compact.
% slanted
\margin{plus simplement, les degrés des diviseurs de $M$ restent fixés \unsure{---|\ldots}}

Donc

\uline{Corollaire} $\underline{\mathrm{Pic}}^n_{X/k}$ est qu. cpct \deleted{\unsure{un}}

Cela équivaut à

\uline{Corollaire} $\underline{\mathrm{Pic}}^n_{X/k}/\mathrm{Pic}^0_{X/k}$ est fini

ou encore

\uline{Corollaire} $\underline{\mathrm{Pic}}^n_{X/k}=\underline{\mathrm{Pic}}^\tau_{X/k}$, le groupe de torsion \deleted{\illegible{1}} de $\text{NérSév}_{X/k}$ est fini

La relation $\underline{\mathrm{Pic}}^{\unsure{n|-|\tau}}_{X/k}=\mathrm{Pic}^\tau_{X/k}$ signifie aussi

\uline{Corollaire} Soit $N=\underline{\mathrm{Pic}}_{X/k}/\underline{\mathrm{Pic}}^\tau_{X/k}$. Alors la forme d'intersection sur $N$ est non singulière.
% archive sheet number: 68

